% !TEX root = ../notes.tex

\documentclass[../notes.tex]{subfiles}

\begin{document}

\section{March 30}
Today, we start talking about automorphic forms.

\subsection{Modular Forms}
We start by reviewing the classical theory of modular forms, which are examples of automorphic forms for the group $\op{SL}_2$. Roughly speaking, a modular form is a symmetric function on the upper-half plane $\mc H$.
\begin{definition}[cusp]
	We set $\mc H\coloneqq\{z\in\CC:\op{Im}z>0\}$ and $\mc H^*\coloneqq\mc H^*\cup\PP^1(\QQ)$. We say that $\PP^1(\QQ)\subseteq\mc H^*$ are the \textit{cusps}.
\end{definition}
\begin{example}
	Note that $\op{SL}_2(\ZZ)\backslash\mc H^*$ is $\op{SL}_2(\ZZ)\backslash\mc H\sqcup\{\infty\}$ because $\op{SL}_2(\ZZ)$
\end{example}
\begin{example}
	For any finite-index subgroup $\Gamma\subseteq\op{SL}_2(\ZZ)$, we see that $\Gamma\backslash\mc H^*=\Gamma\backslash\mc H\sqcup$
\end{example}
\begin{definition}[modular form]
	Fix a congruence subgroup $\Gamma\subseteq\op{SL}_2(\ZZ)$ (meaning that it contains the kernel of the projection $\op{SL}_2(\ZZ)\to\op{SL}_2(\ZZ/N\ZZ)$ for some $N$), and choose an integer $k$ such that $k\ge2$. Then a \textit{modular form of weight $k$ and level $\Gamma$} is a holomorphic function $f\colon\mc H\to\CC$ satisfying the following.
	\begin{itemize}
		\item Symmetry: for all $\begin{bsmallmatrix}
			a & b \\ c & d
		\end{bsmallmatrix}\in\Gamma$ and $z\in\mc H$, we have
		\[f\left(\frac{az+b}{cz+d}\right)=(cz+d)^kf(z).\]
		\item Cusps: the function $f$ is holomorphic at the cusps, meaning that $f$ extends to a holomorphic function on $\mc H^*$.
	\end{itemize}
\end{definition}
Here, being holomorphic at the cusps means that $f$ can be given a power series expansion around each cusp, which extends to a holomorphic function at the cusp. Notably, $\mc H^*$ can be given the structure of a complex manifold, so $f$ can be viewed as a holomorphic function in a punctured disk around a given cusp. Thus, we see that being able to extend $f$ to the puncture amounts to requiring that $f$ is bounded in an open neighborhood of the cusp.
\begin{example} \label{ex:bounded-at-infinity}
	Let's explain what it means for $f$ to be holomorphic at $\infty\in\PP^1(\QQ)$. Here, note that $f(z)=f(z+N)$ for some sufficiently large positive integer $N$. Then $f$ descends to a function on $\CC/N\ZZ$, which is isomorphic to $\CC^\times$ via the exponential map $z\mapsto\exp(2\pi iz/N)$. Thus, $f$ can be viewed as a holomorphic function on $\CC^\times$ in terms of the coordinate $q\coloneqq\exp(2\pi iz/N)$, and we are requiring that $f$ should be bounded as $q\to0$, which corresponds to $\op{Im}z\to\infty$.
\end{example}
Once we have extended $f$ to the cusps, it will turn out to be interesting to just look at those modular forms which vanish at the cusps. Indeed, once we have enough modular forms to be non-vanishing at all the cusps (in a linearly independent way), one can always modify $f$ to vanish at the cusps anyway.
\begin{definition}[cusp form]
	A \textit{cusp form} is a modular form which vanishes at the cusps.
\end{definition}
Let's try to view this definition from the perspective of representation theory. For simplicity, we take $k=2$.
\begin{lemma}
	Modular forms of weight $2$ and level $\Gamma$ correspond to $\Gamma$-invariant meromorphic $1$-forms on $\mc H^*$ which are holomorphic everywhere, except possibly simple poles at the cusps.
\end{lemma}
\begin{proof}
	To start, we note that $\mc H$ is an analytic open subset of $\PP^1(\CC)$, and the transformation law for a modular form $f$ corresponds to making $f(z)\,dz$ into a regular $1$-form on $\mc H$ which is $\Gamma$-invariant: indeed,
	\[f\left(\frac{az+b}{cz+d}\right)d\left(\frac{az+b}{cz+d}\right)=f\left(\frac{az+b}{cz+d}\right)(cz+d)^{-2}\,dz,\]
	which is $f(z)\,dz$ by hypothesis on $f$.
	
	It remains to codify the condition that $f$ is holomorphic at the cusps. Well, by translation, we may as well move the cups to $\infty$, where \Cref{ex:bounded-at-infinity} tells us that we are requiring that $f(z)$ admits an expansion $\varphi(q)$ via the variable $q=\exp(2\pi iz)$. Then $f(z)\,dz=\frac1{2\pi i}\varphi(q)\,\frac{dq}q$, so the condition that $\varphi$ is holomorphic at $q=0$ corresponds to $f(z)\,dz$ admitting at worst a simple pole at $q=0$. Translating this around all the cusps completes the proof of the lemma.
\end{proof}
\begin{remark}
	The same argument shows that cusp forms correspond to $\Gamma$-invariant holomorphic $1$-forms on $\mc H^*$.
\end{remark}

\subsection{Double Quotients}
To continue, we note that $\mc H$ comes from representation theory: $\op{SL}_2(\RR)$ acts transitively on $\mc H$ by fractional linear transformations, and the stabilizer of $i$ is exactly $\op{SO}_2(\RR)$. Thus, $\mc H$ can be viewed as the quotient $\op{SL}_2(\RR)/\op{SO}_2(\RR)$, and $\Gamma\backslash\mc H$ is a double quotient $\Gamma\backslash\op{SL}_2(\RR)/\op{SO}_2(\RR)$.

It will be desirable to write this double quotient in the form
\[\op{SL}_2(\QQ)\backslash\op{SL}_2(\AA_\QQ)/K,\]
where $K\subseteq\op{SL}_2(\AA_\QQ)$ is some compact subgroup. To explain how this works, let's pass to a more general setting. Fix a reductive group $G$ over a number field $F$. For each infinite place $v$, let $K_v\subseteq G(F_v)$ be a maximal compact subgroup; and for each finite place $v$, we let $K_v\subseteq G(F_v)$ be an open compact subgroup. We further require that $K_v=G(\OO_v)$ for all but finitely many $v$.
\begin{remark}
	Note that the condition $K_v=G(\OO_v)$ for all but finitely many $v$ makes sense because $G$ admits an integral model over all but finitely many $v$, and the group $G(\OO_v)$ is the same between such integral models for all but finitely many $v$.
\end{remark}
We further define $K_f\coloneqq\prod_{v<\infty}K_v$ and $\Gamma\coloneqq G(\QQ)\cap K_f$.
% for which $G$ admits an integral model over $\OO_v$ (which is true for all but finitely many $v$), let $K_v\coloneqq G(\OO_v)$ (which does not depend on the choice of integral model for all but finitely many $v$).
\begin{example}
	If $G=\op{SL}_{2,\QQ}$, then $K_f=\op{SL}_2(\widehat{\ZZ})$, 
\end{example}
\begin{theorem} \label{thm:get-strong-approx}
	Fix notation as above. For semisimple, simply connected, split groups $G$ over $\QQ$, the canonical map
	\[\Gamma\backslash G(\RR)\to G(\QQ)\backslash G(\AA_\QQ)/K_f\]
	is a bijection.
\end{theorem}
\begin{remark}
	One can weaken the hypothesis that $G$ is split, but one needs something. For example, if $G(\RR)$ is compact, then the source is compact while the target may not be, so our map cannot be surjective. As an example of such a $G$, we may let $F$ be a quadratic imaginary field over $\QQ$ and choose a Hermitian space $V$ over $F$ which is not definite (once base-changed to $\RR$), and then set $G$ to be the corresponding special unitary group $\op{SU}(V)$.
\end{remark}
It is useful to have language for the conclusion of \Cref{thm:get-strong-approx}.
\begin{definition}[strong approximation]
	Fix a reductive group $G$ over a number field $F$. Then $G$ satisfies \textit{strong approximation} for a finite set of place $v$ if and only if $G(F)\cdot G(F_v)$ is dense in $G(\AA_F)$.
\end{definition}
\begin{remark}
	It turns out that the conclusion of \Cref{thm:get-strong-approx} is equivalent to satisfying strong approximation at the infinite place.
\end{remark}
\begin{remark}
	In general, $G$ satisfies strong approximation at the infinite place if $G$ is semisimple, simply connected, and $G$ is an almost-direct product of $\QQ$-simple groups $G_i$ for which $G_i(\RR)$ is non-compact. There is a similar statement for other places.
\end{remark}
We have just spent some time attempting to explain that strong approximation is a fairly common property of groups, but there are counterexamples.
\begin{example}
	Consider the group $G=\mathbb G_m$ over a number field $F$, and take $K_v=\mathbb G_m(\OO_v)$ for all finite $v$. (Note that $G$ is not semisimple!) Then there is a short exact sequence
	\[1\to\OO_F^\times\backslash F_\infty^\times\to F^\times\backslash\AA_F^\times/K_f\to\op{Cl}(F)\to1.\]
	Thus, in general, $\mathbb G_m(F)\cdot \mathbb G_m(F_\infty)$ cannot be dense in $\mathbb G_m(\AA_F)$ because it generates a finite-index subgroup.
\end{example}
The moral of all of this is that $\op{SL}_2$ does satisfy strong approximation (one can also check this directly), so the natural map
\[(\op{SL}_2(\QQ)\cap K_f)\backslash\op{SL}_2(\RR)\to\op{SL}_2(\QQ)\backslash\op{SL}_2(\AA_\QQ)/K_f\]
is a bijection. Thus, our modular forms can be viewed as meromorphic $1$-forms on this double quotient.

\subsection{Automorphic Forms}
We are now ready to generalize modular forms to automorphic forms. This will require a growth condition, for which we need a notion of height.
\begin{definition}[height]
	Fix an embedded affine variety $\iota\colon X\into\AA^n$ over a number field $F$. Then we define the \textit{height} $\norm\cdot_\iota\colon X(\AA_F)\to\RR^+$ to be given by
	\[\norm x_\iota\coloneqq\prod_v\max\{\left|\iota(x)_1\right|_v,\ldots,\left|\iota(x)_n\right|_v,1\}.\]
\end{definition}
\begin{remark}
	By considering the product embedding, we see that any two heights $\norm\cdot_1$ and $\norm\cdot_2$ admit some large $N$ for which $\norm\cdot_1$ is dominated by $\norm\cdot_2^N$.
\end{remark}
\begin{definition}[automorphic form]
	Fix a reductive group $G$ over a number field $F$. An \textit{automorphic form} $\varphi$ for $G$ is a function $\varphi\colon G(\QQ)\backslash G(\AA_F)\to\CC$ satisfying the following.
	\begin{itemize}
		\item Smooth: for any $x\in G(\AA_F)$, the function $G(F_\infty)\to\CC$ given by $g\mapsto\varphi(xg)$ is smooth.
		\item Smooth: for any $x\in G(\AA_F)$, the function $G(\AA_F^\infty)\to\CC$ given by $g\mapsto\varphi(xg)$ is locally constant.
		\item Finiteness: the function $\varphi$ is $K_\infty$-finite.
		\item Finiteness: for each infinite place $v$, smoothness implies that the Lie algebra $\mf g_v$ of $G(F_v)$ may act by derivatives. Then $\varphi$ must be $Z\mf g_v$-finite.
		\item The function $\varphi$ is fixed by some open compact subgroup of $G(\AA_{F}^\infty)$.
		\item Moderate growth: for any height $\norm\cdot$ on $G$, there are constants $C$ and $r$ for which $\left|\varphi(g)\right|\le C\norm g^r$ for all $g\in G(\AA_F)$.
	\end{itemize}
	We let $\mc A(G)$ denote the space of all automorphic forms.
\end{definition}
Let's see how modular forms give rise to automorphic forms.
\begin{example}
	Let's explain this growth condition for $G=\op{SL}_2(\RR)$. Then $\Gamma\backslash\op{SL}_2(\RR)$ can be covered by matrices of the form
	\[\begin{bmatrix}
		1 & x \\ & 1
	\end{bmatrix}\begin{bmatrix}
		a \\ & a^{-1}
	\end{bmatrix}\]
	where $x$ is found in some finite interval, and $a$ is lower-bounded. (One can show this directly for the congruence subgroups.) By the Phragm\'en--Lindel\"of principle, it turns out that $\varphi$ admitting moderate growth simply means that $\varphi(x+iy)$ has at most polynomial growth as $y\to\infty$.
\end{example}
\begin{example}
	More specifically, if $\varphi$ is a holomorphic function on $\Gamma\backslash\op{SL}_2(\RR)$, then the moderate growth condition means that the $q$-series expansion cannot feature any negative powers of $q=e^{2\pi iz/N}$ because the function $y\mapsto e^{2\pi y/N}$ grows faster than any polynomial.
\end{example}
By construction, we see that $\mc A(G)$ is a $(\mf g,K_\infty)$-module. (One needs to check that derivatives of functions of moderate growth continue to have moderate growth.) Similarly, it is automatically a smooth representation of the totally disconnected group $G(\AA^\infty_F)$, and these two actions commute. Thus, $\mc A$ is almost a representation of $G(\AA_F)$!

We would like for $\mc A$ to be (suitably) admissible, but this is not the case. Here is our target definition, and we will try to find admissible parts of $\mc A$.
\begin{definition}[admissible]
	Fix a module $M$ with commuting actions of $(\mf g,K_\infty)$ and a smooth action of $G(\AA_F^\infty)$. Then $M$ is \textit{admissible} if and only if each irreducible representation of $K_\infty\times K_f$ appears with finite multiplicity.
\end{definition}
% In fact, this representation is admissible, for a suitable notion of admissible: each irreducible representation of $K_\infty\times K_f$ appears with finite multiplicity.
\begin{theorem}[Harish-Chandra] \label{thm:automorphic-get-admissible}
	Fix an ideal $J\subseteq Z\mf g$ of finite codimension. Then the subrepresentation $\mc A(J)$ of $\mc A$ annihilated by $J$ is admissible.
\end{theorem}
\begin{example}
	If we further take $K_f$-invariants, then $\mc A(J)^{K_f}$ is an admissible $(\mf g,K_\infty)$-module.
\end{example}
We close class with the main definition of the course.
\begin{definition}[automorphic representation]
	Fix a reductive group $G$ over a number field $F$. An \textit{automorphic representation} of $G$ is an irreducible admissible representation of $(\mf g,K_\infty)$ and $G(\AA_F^\infty)$ (with commuting actions) which appears as a subquotient of $\mc A$.
\end{definition}
\begin{remark}
	Note that automorphic representations admit an infinitesimal character, so it comes from some $\mc A(J)$. It follows that automorphic representations are admissible by \Cref{thm:automorphic-get-admissible}.
\end{remark}

\end{document}